\usepackage{fontspec}
\usepackage{polyglossia}
\setmainlanguage{spanish}
\usepackage{amsmath}
\usepackage{amssymb}
\usepackage{xcolor}
\usepackage[colorlinks=true,linkcolor=azulLC,urlcolor=azulLC,
            allbordercolors={1 1 1}]{hyperref}
\usepackage{listings}
\usepackage{geometry}
\geometry{margin=2.5cm}

% ── Tema oscuro estilo VS Code (Dark+) ──
% Fondo de página, texto y sintaxis imitan el editor VS Code para lectura nocturna.
\definecolor{fondoPagina}{HTML}{1E1E1E}   % editor background
\definecolor{textoPrincipal}{HTML}{D4D4D4}% editor foreground
\definecolor{azulLC}{HTML}{3794FF}        % enlaces (textLink de VS Code)
\definecolor{fondoCodigo}{HTML}{252526}   % sidebar/panel background
\definecolor{bordeCodigo}{HTML}{3C3C3C}   % borde sutil del bloque de código
\definecolor{comentario}{HTML}{6A9955}    % comentarios (verde)
\definecolor{palabraClave}{HTML}{569CD6}  % keywords (azul)
\definecolor{cadena}{HTML}{CE9178}        % strings (salmón)
\definecolor{numeroLinea}{HTML}{858585}   % números de línea (gris)

% Fondo oscuro global + texto claro por defecto.
% Redefinir \normalcolor es CLAVE: muchos internos de LaTeX (encabezados,
% títulos, reglas de tabular/\hline, líneas guía del índice) llaman
% \normalcolor y volverían a NEGRO. Al apuntarlo al color claro, nada queda
% en negro sobre el fondo oscuro.
\pagecolor{fondoPagina}
\AtBeginDocument{%
  \color{textoPrincipal}%
  \renewcommand{\normalcolor}{\color{textoPrincipal}}%
}

% Estilo de listado C++
\lstdefinestyle{cpp}{
  language=C++,
  backgroundcolor=\color{fondoCodigo},
  basicstyle=\ttfamily\small\color{textoPrincipal},
  keywordstyle=\color{palabraClave}\bfseries,
  commentstyle=\color{comentario}\itshape,
  stringstyle=\color{cadena},
  numbers=left,
  numberstyle=\tiny\color{numeroLinea},
  numbersep=8pt,
  showstringspaces=false,
  breaklines=true,
  frame=single,
  rulecolor=\color{bordeCodigo},
  tabsize=2,
  extendedchars=true,
  literate=%
    {á}{{\'a}}1 {é}{{\'e}}1 {í}{{\'i}}1 {ó}{{\'o}}1 {ú}{{\'u}}1
    {ñ}{{\~n}}1 {Á}{{\'A}}1 {É}{{\'E}}1 {Í}{{\'I}}1 {Ó}{{\'O}}1
    {Ú}{{\'U}}1 {Ñ}{{\~N}}1 {¿}{{?`}}1 {¡}{{!`}}1
}
\lstset{style=cpp}
